% Fragment -- inclus de main.tex prin \input
\clearpage
\setpartlabel{themeJs}{\faBolt}{Faza 3: JavaScript}
\setcounter{section}{0}
\phantomsection
\addcontentsline{toc}{part}{\protect\textbf{Faza 3 \textendash{} JavaScript}}
\handout{Faza 3 \textendash{} JavaScript}{Variabile. Functii. Pattern-ul SELECT-LISTEN-CHANGE. Manipulare DOM.}

JavaScript este limbajul de programare standardizat ECMAScript (ECMA-262), interpretat nativ de toate browserele moderne. Spre deosebire de HTML (structura) si CSS (aspect), JavaScript adauga \textbf{comportament}: raspunde la actiunile utilizatorului, modifica continutul paginii dinamic, comunica cu servere prin retea.

Capitolul prezinta sintaxa de baza, modelul de programare orientat-eveniment specific paginilor web statice, si interactiunea cu arborele DOM (\emph{Document Object Model}).

\bigskip

\begin{outcomes}
Dupa capitolul asta vei \textcommabelow{s}ti:
\begin{itemize}[itemsep=1pt]
  \item Declara variabile cu \code{const} \textcommabelow{s}i \code{let}, recunoscand cand fiecare se aplica.
  \item Aplica pattern-ul \textbf{SELECT $\to$ LISTEN $\to$ CHANGE} pentru orice interac\textcommabelow{t}iune.
  \item Manipuleaza DOM-ul cu \code{querySelector}, \code{classList}, \code{addEventListener}.
  \item Intercepteaza submit-ul unui formular \textcommabelow{s}i extrage valorile cu \code{FormData}.
  \item Folose\textcommabelow{s}te \code{IntersectionObserver} pentru anima\textcommabelow{t}ii \textcommabelow{s}i lazy loading legate de scroll.
  \item Diagnostiheaza erorile JavaScript folosind console-ul browserului.
  \item \faRocket\ \emph{Bonus self-study:} debounce/throttle, event delegation, \code{requestAnimationFrame}, \code{fetch} cu async/await, metode pe array (\code{map}/\code{filter}), \code{localStorage}.
\end{itemize}
\end{outcomes}

\begin{whymatters}
HTML \textcommabelow{s}i CSS construiesc o pagina statica. JavaScript transforma pagina intr-o aplica\textcommabelow{t}ie: poate reac\textcommabelow{t}iona la click-uri, poate schimba con\textcommabelow{t}inutul fara reincarcare, poate trimite cereri catre servere. Pattern-ul SELECT-LISTEN-CHANGE prezentat aici acopera 80\% din scenariile de interactivitate \textcommabelow{s}i este aceea\textcommabelow{s}i abstrac\textcommabelow{t}ie pe care React, Vue \textcommabelow{s}i Svelte o ambala in API-uri declarative.
\end{whymatters}

\begin{nota}[Pre-rechizite]
Capitolul presupune in\textcommabelow{t}elegerea structurii HTML (tag-uri, atribute) \textcommabelow{s}i a selectorilor CSS (clase, id-uri). Faze 1 \textcommabelow{s}i 2.
\end{nota}

% =========================================================
\section{Conectarea fisierului JavaScript la HTML}

Includerea unui fisier JavaScript se realizeaza prin elementul \code{<script>}, plasat in \code{<head>} cu atributul \code{defer}:

\begin{lstlisting}[style=html]
<script src="script.js" defer></script>
\end{lstlisting}

\subsection{\texttt{defer} vs \texttt{async} vs nimic}

\begin{itemize}
  \item \textbf{Fara atribut.} Browserul opreste parsarea HTML-ului, descarca si ruleaza scriptul, apoi continua. Blocheaza randarea \textendash{} antipattern.
  \item \code{async} \textendash{} descarcare in paralel cu parsarea, dar executie imediat ce scriptul este disponibil (poate intrerupe parsarea). Util doar pentru scripturi independente (analytics).
  \item \code{defer} \textendash{} descarcare in paralel, executie \textbf{dupa} parsarea HTML-ului. Garantat executat in ordinea declararii. Optiunea recomandata pentru codul propriu.
\end{itemize}

\begin{atentie}
Fara \code{defer}, scriptul ruleaza inainte de parsarea \code{<body>}: \code{document.querySelector('.cta')} returneaza \code{null}. Cu \code{defer}, intregul DOM exista cand ruleaza scriptul.
\end{atentie}

% =========================================================
\section{Declararea variabilelor}

ECMAScript 2015 introduce doua cuvinte cheie pentru declararea variabilelor: \code{const} (valoare imuabila) si \code{let} (valoare mutabila). Cuvantul cheie \code{var} este considerat depasit din cauza regulilor de scope.

\begin{lstlisting}[style=js, caption={Reguli de declarare}]
const PI = 3.14;            // valoare imutabila
let counter = 0;            // valoare mutabila

counter = counter + 1;      // permis
// PI = 3.15;               // TypeError la rulare
\end{lstlisting}

\subsection{Diferenta dintre \texttt{var} si \texttt{let} / \texttt{const}}

\code{var} are \emph{function scope}: o variabila declarata in interiorul unui \code{if} sau \code{for} este accesibila in tot blocul functiei. \code{let} si \code{const} au \emph{block scope}: vizibile doar in interiorul \code{\{...\}} in care au fost declarate.

\begin{dano}[Asa nu \textendash{} \texttt{var}, scope confuz]
\begin{lstlisting}[style=js, numbers=none]
function exemplu() {
  if (true) {
    var x = 10;
  }
  console.log(x);  // 10 -- scapa din if!
}
\end{lstlisting}
\end{dano}

\begin{dado}[Asa da \textendash{} \texttt{let}, scope previzibil]
\begin{lstlisting}[style=js, numbers=none]
function exemplu() {
  if (true) {
    let x = 10;
  }
  console.log(x);  // ReferenceError -- nu exista in afara if-ului
}
\end{lstlisting}
\end{dado}

\textbf{Conventia practica:} \code{const} este valoarea implicita; \code{let} se foloseste doar cand reasignarea este necesara; \code{var} se evita complet in cod nou.

% =========================================================
\section{Tipuri primitive si compuse}

\begin{lstlisting}[style=js]
const nume      = "Voluntar";              // string
const varsta    = 21;                   // number (intregi si zecimale)
const eStudent  = true;                 // boolean
const limbaje   = ["HTML", "CSS"];      // array (lista ordonata)
const persoana  = {                     // object (pereche cheie-valoare)
  nume:   "Voluntar",
  varsta: 21,
};

console.log(persoana.nume);             // "Voluntar"
console.log(limbaje[0]);                // "HTML"
console.log(typeof varsta);             // "number"
\end{lstlisting}

JavaScript este un limbaj cu \emph{tipare dinamica}: tipul unei variabile este determinat de valoarea atribuita si poate fi modificat in timpul executiei.

\subsection{Verificarea egalitatii: \texttt{===} vs \texttt{==}}

Operatorul \code{==} efectueaza \emph{coercitie de tip}: converteste operanzii inainte de comparatie. Rezultatul este des contraintuitiv. Operatorul \code{===} compara strict valoare si tip.

\begin{dano}[Asa nu \textendash{} \texttt{==}, coercitie surprinzatoare]
\begin{lstlisting}[style=js, numbers=none]
0 == "0"             // true
0 == ""              // true
"" == false          // true
null == undefined    // true
[] == false          // true
\end{lstlisting}
\end{dano}

\begin{dado}[Asa da \textendash{} \texttt{===} mereu]
\begin{lstlisting}[style=js, numbers=none]
0 === "0"            // false (number vs string)
0 === ""             // false
null === undefined   // false
\end{lstlisting}
\end{dado}

% =========================================================
\section{Functii: declaratie, expresie, arrow}

JavaScript ofera trei sintaxe pentru definirea unei functii. Cele trei comportamente difera la nivel de \code{this} si hoisting, dar pentru cazul uzual (handler de eveniment) sunt interschimbabile.

\begin{lstlisting}[style=js, caption={Trei moduri de a scrie aceeasi functie}]
// 1. Declaratie de functie (function declaration)
function saluta(nume) {
  return `Salut, ${nume}!`;
}

// 2. Expresie de functie (function expression)
const saluta2 = function(nume) {
  return `Salut, ${nume}!`;
};

// 3. Arrow function (sintaxa concisa)
const saluta3 = (nume) => `Salut, ${nume}!`;

// Apel: identic in toate cazurile
console.log(saluta("Voluntar"));
\end{lstlisting}

\textbf{Conventia practica:} arrow functions pentru handlere si callbacks; declaratii pentru functii cu nume reutilizate, definite la nivelul fisierului.

% =========================================================
\section{Pattern-ul SELECT $\to$ LISTEN $\to$ CHANGE}

Majoritatea interactiunilor pe paginile statice respecta o schema in trei pasi: selectarea elementului DOM, atasarea unui ascultator de eveniment, modificarea elementului in raspuns la eveniment.

\begin{lstlisting}[style=js, caption={Toggle pentru tema dark/light}]
// 1. SELECT
const btn = document.querySelector('.toggle-theme');

// 2. LISTEN
btn.addEventListener('click', () => {

  // 3. CHANGE
  document.body.classList.toggle('dark');

});
\end{lstlisting}

\begin{ponturi}
Pattern-ul descris reprezinta abstractizarea fundamentala pe care framework-urile moderne (React, Vue, Svelte, Angular) o ascund in spatele propriilor API-uri declarative. Intelegerea variantei \emph{vanilla} faciliteaza tranzitia ulterioara catre framework-uri.
\end{ponturi}

% =========================================================
\section{Manipularea DOM-ului}

API-ul \code{document} expune metode pentru selectia si modificarea nodurilor:

\begin{lstlisting}[style=js, caption={Operatii uzuale}]
// SELECT -- un singur element
const el = document.querySelector('.cta');           // primul .cta
const id = document.getElementById('contact');       // alternativa pt id

// SELECT -- mai multe elemente
const linkuri = document.querySelectorAll('nav a');
linkuri.forEach((link) => {
  link.addEventListener('click', () => console.log('click'));
});

// CHANGE -- continut
el.textContent = 'Salut!';                  // text simplu (sigur)
el.innerHTML  = '<strong>Salut!</strong>';  // HTML (atentie XSS)

// CHANGE -- atribute
el.setAttribute('aria-expanded', 'true');
el.classList.add('active');
el.classList.remove('hidden');
el.classList.toggle('open');

// CHANGE -- creare element nou
const nou = document.createElement('li');
nou.textContent = 'Element nou';
document.querySelector('ul').appendChild(nou);
\end{lstlisting}

\begin{dano}[Asa nu \textendash{} querySelector la fiecare apel]
\begin{lstlisting}[style=js, numbers=none]
btn.addEventListener('click', () => {
  document.querySelector('.modal').classList.add('open');
  document.querySelector('.modal').focus();
  document.querySelector('.modal').scrollIntoView();
});
\end{lstlisting}
\textbf{Problema.} Trei traversari ale DOM-ului pentru aceeasi referinta.
\end{dano}

\begin{dado}[Asa da \textendash{} cache local]
\begin{lstlisting}[style=js, numbers=none]
const modal = document.querySelector('.modal');

btn.addEventListener('click', () => {
  modal.classList.add('open');
  modal.focus();
  modal.scrollIntoView();
});
\end{lstlisting}
\end{dado}

% =========================================================
\section{Procesarea formularelor}

Submit-ul nativ al unui formular trimite o cerere HTTP catre URL-ul declarat in \code{action} si reincarca pagina. In majoritatea cazurilor, fluxul dorit este interceptarea evenimentului si manipularea datelor in JavaScript:

\begin{lstlisting}[style=js, caption={Captura datelor de formular fara reincarcarea paginii}]
const form = document.querySelector('#contact-form');

form.addEventListener('submit', (e) => {
  e.preventDefault();                                   // blocheaza reincarcarea

  const data = Object.fromEntries(new FormData(form));  // {nume, email, mesaj}
  console.log('Date primite:', data);

  form.reset();                                         // goleste campurile
  alert(`Multumesc, ${data.nume}!`);
});
\end{lstlisting}

Fluxul: \code{e.preventDefault()} suspenda comportamentul implicit; \code{new FormData(form)} colecteaza valorile dupa atributul \code{name}; \code{Object.fromEntries(...)} le converteste intr-un obiect uzual; \code{form.reset()} reseteaza valorile in interfata.

\begin{ponturi}
Pentru transmiterea efectiva a datelor catre un endpoint extern (fara backend propriu), serviciul \href{https://formspree.io}{Formspree} ofera un endpoint gratuit. Configurarea consta intr-o singura modificare a atributului \code{action}.
\end{ponturi}

% =========================================================
\section{Debugging si console}

Console-ul browserului (F12 $\to$ tab \emph{Console}) este principalul instrument de diagnostic. Erorile JavaScript apar automat aici cu mesaj, fisier si linie.

\begin{lstlisting}[style=js, caption={Trei moduri de a inspecta valorile}]
// Cel mai uzual: log valoare cu eticheta
console.log('Date primite:', data);

// Tabel pentru obiecte / array-uri (formatare imbunatatita)
console.table(linkuri);

// Pauza in executie -- DevTools deschis
debugger;
\end{lstlisting}

\textbf{Sfaturi practice.} Inainte de a cere ajutor, citeste mesajul de eroare din console. Numele functiei si numarul liniei sunt aproape intotdeauna corecte. \code{console.log(typeof x)} verifica tipul cand comportamentul este surprinzator.

% =========================================================
\section{IntersectionObserver \textendash{} anima\textcommabelow{t}ii la scroll}

Cea mai uzuala cerere ,,fa o anima\textcommabelow{t}ie cand sectiunea devine vizibila'' (fade-in pe scroll, reveal pe carduri, lazy loading pentru imagini) NU se rezolva cu \code{window.addEventListener('scroll', ...)}. Browserul ofera un API dedicat, mult mai eficient: \code{IntersectionObserver}.

Idea: in loc sa intrebi de 60 de ori pe secunda ,,e vizibila sectiunea?'', tu declari ,,anun\textcommabelow{t}a-ma cand starea de vizibilitate se schimba''. Browserul face restul, fara cost de performance.

\begin{lstlisting}[style=js, caption={Fade-in pe scroll: pattern canonic in 10 linii}]
const observer = new IntersectionObserver((entries) => {
  entries.forEach((entry) => {
    if (entry.isIntersecting) {
      entry.target.classList.add('visible');
      observer.unobserve(entry.target);  // o singura data, nu repeti
    }
  });
}, { threshold: 0.15 });  // declanseaza cand 15% e in viewport

document.querySelectorAll('.reveal').forEach((el) => observer.observe(el));
\end{lstlisting}

Pe CSS, partea sora:

\begin{lstlisting}[style=css, numbers=none]
.reveal { opacity: 0; transform: translateY(20px); transition: all 0.6s ease-out; }
.reveal.visible { opacity: 1; transform: translateY(0); }
\end{lstlisting}

Cum se cite\textcommabelow{s}te API-ul:
\begin{itemize}[itemsep=1pt]
  \item Constructorul: \code{new IntersectionObserver(callback, optiuni)}.
  \item \code{callback(entries)} \textendash{} se ape\-leaza cand starea de vizibilitate se schimba.
  \item \code{entry.isIntersecting} \textendash{} \code{true} cand elementul intra in viewport.
  \item \code{entry.target} \textendash{} elementul DOM (acela\textcommabelow{s}i pe care l-ai dat la \code{observe}).
  \item \code{threshold} \textendash{} procent din element vizibil pentru a declan\textcommabelow{s}a (\code{0} = orice pixel, \code{1} = complet vizibil, \code{0.15} = 15\%).
  \item \code{rootMargin} \textendash{} extinde/contracta viewport-ul (ex. \code{"0px 0px -100px 0px"} declan\textcommabelow{s}eaza cu 100px inainte sa intre real in cadru).
  \item \code{observer.unobserve(el)} \textendash{} opre\textcommabelow{s}te monitorizarea (perfect dupa prima activare).
\end{itemize}

\begin{dano}[Asa nu \textendash{} scroll handler cu \texttt{getBoundingClientRect}]
\begin{lstlisting}[style=js, numbers=none]
window.addEventListener('scroll', () => {
  document.querySelectorAll('.reveal').forEach((el) => {
    if (el.getBoundingClientRect().top < window.innerHeight) {
      el.classList.add('visible');
    }
  });
});
\end{lstlisting}
\textbf{Problema.} Ruleaza la fiecare pixel de scroll (zeci \textendash{} sute de ori pe secunda), forteaza recalculare layout (\emph{layout thrashing}), scade FPS pe mobile.
\end{dano}

\begin{dado}[Asa da \textendash{} IntersectionObserver]
Browserul observa eficient (compus pe thread separat). Callback-ul ruleaza doar cand starea se schimba. Zero impact pe performance.
\end{dado}

\begin{ponturi}[Aplica\textcommabelow{t}ii directe]
\begin{itemize}[itemsep=1pt]
  \item Lazy loading custom pentru con\textcommabelow{t}inut dinamic (cand nu poti folosi \code{loading="lazy"}).
  \item Active state pentru navigare \emph{scrollspy} (eviden\textcommabelow{t}iaza link-ul activ in nav).
  \item Infinite scroll \textendash{} cand userul ajunge aproape de jos, incarci urmatoarea pagina.
  \item Analytics ,,a vazut sec\textcommabelow{t}iunea X'' (tracking eveniment doar cand sec\textcommabelow{t}iunea devine vizibila).
\end{itemize}
\end{ponturi}

% =========================================================
\section{Tratarea erorilor (\texttt{try}/\texttt{catch})}

Codul care poate gen\-era erori la rulare (parsare JSON, apeluri \code{fetch}, acces la API-uri externe) trebuie incadrat intr-un bloc \code{try}/\code{catch} pentru ca o eroare sa nu opreasca executia restului paginii:

\begin{lstlisting}[style=js, caption={Parsare JSON cu fallback}]
const raw = localStorage.getItem('preferinte');

try {
  const config = JSON.parse(raw);
  console.log('Preferinte incarcate:', config);
} catch (err) {
  console.warn('JSON invalid in localStorage; folosesc valori implicite');
  // continuare cu valori default
}
\end{lstlisting}

Cand eroarea este previzibila si recuperabila (utilizatorul poate continua), \code{try}/\code{catch} este abordarea corecta. Pentru erori neprevazute, evenimentul global \code{window.addEventListener('error', ...)} permite logarea catre un serviciu de monitoring.

% =========================================================
\section{Erori frecvente}

\begin{itemize}
  \item \textbf{\code{Cannot read property '...' of null}.} Cauza: \code{querySelector} a returnat \code{null} si codul a incercat sa acceseze o proprietate. Diagnostic: \code{console.log(el)} inainte de a folosi.
  \item \textbf{\code{addEventListener is not a function}.} Aceeasi cauza: rezultatul \code{querySelector} este \code{null}.
  \item \textbf{Hoisting.} Variabilele \code{let} si \code{const} declarate intr-un bloc nu sunt accesibile inainte de declaratie. Accesul anticipat genereaza \emph{ReferenceError: Cannot access before initialization}.
  \item \textbf{Comparatie cu \code{undefined}.} Foloseste \code{x === undefined}, nu \code{!x} \textendash{} ultimul prinde si \code{0}, \code{""}, \code{false}, \code{null}.
\end{itemize}

% =========================================================
\begin{bonusbox}
\textbf{Acest box e optional.} Daca esti la inceput, treci direct la \emph{Verificare cuno\textcommabelow{s}tin\textcommabelow{t}e}. Mai jos sunt 8 pattern-uri pe care le folose\textcommabelow{s}ti zilnic in proiecte reale.

\medskip

\textbf{\faClock\ Debounce \textcommabelow{s}i throttle \textendash{} st\u apane\textcommabelow{s}te evenimentele frecvente}

Evenimentele \code{scroll}, \code{resize}, \code{input}, \code{mousemove} se declan\textcommabelow{s}eaza de zeci \textendash{} sute de ori pe secunda. Daca handler-ul face ceva costisitor (fetch, recalcul layout), pagina ,,inghea\textcommabelow{t}a''.

\emph{Debounce} \textendash{} amana executia pana cand evenimentul nu mai apare \emph{ms} milisecunde. Bun pentru autocomplete (a\textcommabelow{s}tepta sa termini de scris).

\emph{Throttle} \textendash{} executa cel mult o data la fiecare \emph{ms}. Bun pentru scroll, resize.

\begin{lstlisting}[style=js, caption={Utilitar reutilizabil}]
function debounce(fn, ms = 300) {
  let timer;
  return function(...args) {
    clearTimeout(timer);
    timer = setTimeout(() => fn.apply(this, args), ms);
  };
}

const search = document.querySelector('#search');
search.addEventListener('input', debounce((e) => {
  console.log('Cautare:', e.target.value);  // ruleaza dupa 300ms de tacere
}, 300));
\end{lstlisting}

\begin{lstlisting}[style=js, caption={Buton ,,scroll to top'' care apare dupa 400px}]
function throttle(fn, ms = 100) {
  let last = 0;
  return function(...args) {
    const now = Date.now();
    if (now - last >= ms) { last = now; fn.apply(this, args); }
  };
}

const btnTop = document.querySelector('.to-top');
window.addEventListener('scroll', throttle(() => {
  btnTop.classList.toggle('visible', window.scrollY > 400);
}, 100));
\end{lstlisting}

\medskip

\textbf{\faSitemap\ Event delegation \textendash{} un singur listener pentru N elemente}

Daca ai 50 de butoane \code{.delete-btn}, NU pune 50 de \code{addEventListener}. Pune unul pe parinte si folose\textcommabelow{s}te \code{event.target.closest()} pentru a verifica ce a fost clicat. Bonus: func\textcommabelow{t}ioneaza \textcommabelow{s}i pentru elemente adaugate dinamic dupa atasare.

\begin{lstlisting}[style=js, caption={Lista de task-uri cu \u stergere}]
const list = document.querySelector('#task-list');

list.addEventListener('click', (e) => {
  const deleteBtn = e.target.closest('.delete-btn');
  if (!deleteBtn) return;                          // click in alta parte
  deleteBtn.closest('li').remove();
});
\end{lstlisting}

Cum func\textcommabelow{t}ioneaza: evenimentul \emph{bubbles} (urca de la \code{target} catre parinti). Listener-ul pe parinte vede toate click-urile copiilor.

\medskip

\textbf{\faCog\ Op\textcommabelow{t}iuni la \texttt{addEventListener}}

Parametrul al treilea poate fi un obiect cu op\textcommabelow{t}iuni:

\begin{itemize}[itemsep=1pt]
  \item \code{\{ once: true \}} \textendash{} executa o singura data, apoi se auto-elimina. Util pentru splash/intro animation.
  \item \code{\{ passive: true \}} \textendash{} promiti browserului ca NU vei face \code{preventDefault()}. Pentru \code{scroll}/\code{touchmove} mai ales \textendash{} cre\textcommabelow{s}te FPS pe mobile.
  \item \code{\{ capture: true \}} \textendash{} prinde evenimentul in faza de \emph{capture} (de sus in jos), nu \emph{bubble} (rar folosit).
\end{itemize}

\begin{lstlisting}[style=js, numbers=none]
window.addEventListener('scroll', onScroll, { passive: true });
btn.addEventListener('click', showWelcome, { once: true });
\end{lstlisting}

\medskip

\textbf{\faPlay\ \texttt{requestAnimationFrame} \textendash{} anima\textcommabelow{t}ii fluide din JS}

Cand vrei sa animezi ceva din JavaScript (nu doar prin CSS), NU folosi \code{setTimeout(..., 16)}. Folose\textcommabelow{s}te \code{requestAnimationFrame} \textendash{} se sincronizeaza cu refresh-ul ecranului (60fps natural, 120fps pe ecrane rapide).

\begin{lstlisting}[style=js, caption={Counter animat de la 0 la o valoare \textcommabelow{t}int\u a}]
function animateNumber(el, target, duration = 1000) {
  const start = performance.now();

  function tick(now) {
    const elapsed = now - start;
    const progress = Math.min(elapsed / duration, 1);   // 0..1
    el.textContent = Math.floor(progress * target);
    if (progress < 1) requestAnimationFrame(tick);
  }

  requestAnimationFrame(tick);
}

animateNumber(document.querySelector('.stat-projects'), 42);
\end{lstlisting}

\medskip

\textbf{\faGlobe\ Async / await + fetch \textendash{} apeluri spre API-uri}

Functiile \code{async} permit scrierea codului asincron in stil sincron. \code{await} pauzeaza pana cand un \code{Promise} se rezolva.

\begin{lstlisting}[style=js, caption={Voluntar afi\textcommabelow{s}eaz\u a stats GitHub live pe portofoliu}]
async function loadGitHubStats(username) {
  try {
    const response = await fetch(`https://api.github.com/users/${username}`);
    if (!response.ok) throw new Error(`HTTP ${response.status}`);

    const data = await response.json();

    document.querySelector('.gh-repos').textContent     = data.public_repos;
    document.querySelector('.gh-followers').textContent = data.followers;
    document.querySelector('.gh-avatar').src            = data.avatar_url;
  } catch (err) {
    console.warn('Nu am putut incarca stats GitHub:', err);
  }
}

loadGitHubStats('voluntar-best');
\end{lstlisting}

Trei reguli:
\begin{itemize}[itemsep=1pt]
  \item \code{async} inainte de \code{function} \textendash{} func\textcommabelow{t}ia returneaza un \code{Promise}.
  \item \code{await} func\textcommabelow{t}ioneaza doar in interiorul unei func\textcommabelow{t}ii \code{async}.
  \item Erorile asincrone se prind cu \code{try}/\code{catch} la fel ca cele sincrone.
\end{itemize}

\medskip

\textbf{\faList\ Metode pe array \textendash{} \texttt{map}, \texttt{filter}, \texttt{find}}

In loc de \code{for} cu indice si \code{array.push()}, JavaScript ofera metode care returneaza un array nou (imuabilitate). Mult mai citibile pentru transformari de date.

\begin{lstlisting}[style=js, caption={Lista de proiecte randat\u a din date}]
const proiecte = [
  { titlu: 'Portfolio',   link: '/portfolio',   tag: 'web' },
  { titlu: 'Game of Life', link: '/gol',         tag: 'algo' },
  { titlu: 'Calculator',   link: '/calc',        tag: 'web' },
];

// map -- transforma fiecare element
const html = proiecte
  .filter((p) => p.tag === 'web')                      // doar proiectele web
  .map((p) => `<li><a href="${p.link}">${p.titlu}</a></li>`)
  .join('');

document.querySelector('.projects').innerHTML = html;

// find -- gase\textcommabelow{s}te primul element care satisface
const portfolio = proiecte.find((p) => p.titlu === 'Portfolio');
\end{lstlisting}

\code{reduce} (mai rar in front-end de inceput) acumuleaza un singur rezultat din array (suma, grup\u ari, statistici).

\medskip

\textbf{\faPuzzlePiece\ Destructuring + template literals}

Doua sintaxe pe care le vezi peste tot. Destructuring extrage valori din obiecte/array-uri intr-un singur pas:

\begin{lstlisting}[style=js, numbers=none]
const { nume, varsta } = persoana;       // in loc de persoana.nume, persoana.varsta
const [primul, al_doilea] = limbaje;     // in loc de limbaje[0], limbaje[1]

// Aplicat in parametri de functie
function saluta({ nume, varsta }) {
  return `${nume}, ${varsta} ani`;
}
saluta({ nume: 'Voluntar', varsta: 21 });
\end{lstlisting}

Template literals (backtick) permit interpolare \code{\$\{...\}} \textcommabelow{s}i string-uri multi-linie. Le-ai folosit deja in exemplul GitHub si in HTML-ul generat din \code{map}.

\medskip

\textbf{\faSave\ \texttt{localStorage} \textendash{} salveaza preferin\textcommabelow{t}ele user-ului}

\code{localStorage} este un dictionar persistent intre vizite. Stocheaza doar string-uri \textendash{} pentru obiecte folose\textcommabelow{s}ti \code{JSON.stringify}/\code{JSON.parse}.

\begin{lstlisting}[style=js, caption={Tema dark se p\u astreaz\u a intre vizite}]
const toggle = document.querySelector('.toggle-theme');

// La incarcare: aplica tema salvata
if (localStorage.getItem('theme') === 'dark') {
  document.body.classList.add('dark');
}

// La click: schimba si salveaza
toggle.addEventListener('click', () => {
  document.body.classList.toggle('dark');
  const tema = document.body.classList.contains('dark') ? 'dark' : 'light';
  localStorage.setItem('theme', tema);
});
\end{lstlisting}

\code{sessionStorage} e identic, dar se sterge cand userul inchide tab-ul. Ambele au limita ~5MB per origin.

\medskip

\textbf{\faRoute\ Unde mergi de aici}

\begin{itemize}[itemsep=1pt]
  \item \emph{MDN IntersectionObserver:} \url{https://developer.mozilla.org/en-US/docs/Web/API/Intersection_Observer_API}
  \item \emph{javascript.info \textendash{} Promises, async/await:} \url{https://javascript.info/async}
  \item \emph{web.dev Event handling best practices:} \url{https://web.dev/articles/optimize-input-delay}
  \item \emph{Free public APIs pentru exersat:} \url{https://github.com/public-apis/public-apis}
\end{itemize}
\end{bonusbox}

% =========================================================
\section*{Verificare cuno\textcommabelow{s}tin\textcommabelow{t}e}

\begin{selfcheck}
\begin{enumerate}[itemsep=2pt]
  \item Ce face atributul \code{defer} pe un \code{<script>} \textcommabelow{s}i de ce este necesar?
  \item Care sunt cei trei pa\textcommabelow{s}i ai pattern-ului \textbf{SELECT $\to$ LISTEN $\to$ CHANGE}?
  \item De ce \code{===} este preferabil lui \code{==}? Da\textcommabelow{t}i un exemplu de coerci\textcommabelow{t}ie surprinzatoare.
  \item Ce face \code{e.preventDefault()} apelat intr-un handler de \code{submit}?
  \item \emph{Spot the issue:} \code{const btn = document.querySelector('.cta'); btn.addEventListener(...);} \textendash{} \code{btn} este \code{null}.
  \item De ce este \code{IntersectionObserver} preferabil lui \code{scroll} listener pentru anima\textcommabelow{t}ii la scroll?
  \item Ce face \code{observer.unobserve(el)} \textcommabelow{s}i cand il folose\textcommabelow{s}ti?
\end{enumerate}
\end{selfcheck}

% =========================================================
\section*{Recapitulare}

\begin{recap}[Faza 3 -- JavaScript]
\begin{itemize}[itemsep=1pt]
  \item Includere: \code{<script src="..." defer></script>} in \code{<head>}.
  \item \code{const} este declararea implicita; \code{let} doar cand reasignarea este necesara; nu \code{var}.
  \item Comparatie: exclusiv \code{===}, niciodata \code{==}.
  \item Pattern-cheie: \textbf{SELECT $\to$ LISTEN $\to$ CHANGE} acopera majoritatea interactiunilor.
  \item DOM: \code{querySelector} (un element); \code{querySelectorAll} (NodeList iterabil cu \code{.forEach}).
  \item Form submit: \code{e.preventDefault()} + \code{new FormData(form)} + \code{Object.fromEntries(...)}.
  \item Cache referintele DOM in variabile; nu repeta \code{querySelector} pentru acelasi element.
  \item \code{IntersectionObserver} pentru anima\textcommabelow{t}ii la scroll: zero cost de performance, callback doar cand vizibilitatea se schimba.
  \item Debug: \code{F12 $\to$ Console}, \code{console.log()}, \code{debugger}.
  \item \faRocket\ Bonus: debounce/throttle, event delegation, \code{requestAnimationFrame}, \code{fetch} cu async/await, \code{map}/\code{filter}/\code{find}, destructuring, \code{localStorage}.
\end{itemize}
\end{recap}

\partdivider

% =========================================================
\begin{joc}
\textbf{Platforma:} CodePen (\url{codepen.io}). \quad \textit{URL-ul exact este distribuit de trainer in chat.}

\medskip

\textbf{Activitatea.} Documentul de pornire contine HTML si CSS complet, dar fisierul JavaScript este gol. Sarcina este implementarea pattern-ului SELECT-LISTEN-CHANGE: la apasarea butonului prezent in pagina, clasa \code{dark} se aplica si se scoate alternativ pe elementul \code{<body>}.

\medskip

\textbf{Solutia asteptata:}

\begin{lstlisting}[style=js, numbers=none]
const btn = document.querySelector('button');
btn.addEventListener('click', () => {
  document.body.classList.toggle('dark');
});
\end{lstlisting}

\textbf{Durata estimata:} 5 minute.
\end{joc}

% =========================================================
\section*{\faGraduationCap\ \ Resurse pentru aprofundare}

\begin{itemize}[itemsep=2pt]
  \item Kantor I. \emph{The Modern JavaScript Tutorial}. \url{https://javascript.info}
  \item Mozilla Developer Network. \emph{JavaScript reference}. \\\url{https://developer.mozilla.org/en-US/docs/Web/JavaScript}
  \item The Odin Project. \emph{JavaScript Foundations}. \\\url{https://www.theodinproject.com/paths/foundations/courses/foundations\#javascript-basics}
  \item Mozilla Developer Network. \emph{Document Object Model}. \\\url{https://developer.mozilla.org/en-US/docs/Web/API/Document_Object_Model}
  \item Google. \emph{Learn JavaScript}. \url{https://web.dev/learn/javascript}
\end{itemize}

\bigskip

\bridge{Capitolul 4 (Deploy) muta site-ul de pe ma\textcommabelow{s}ina locala pe un URL public. Toate fazele 1\textendash{}3 sunt locale; abia dupa deploy pagina exista in lume.}
